% % \clearpage


\section{Related Work}
We review the related methods from two aspects, including the visual dubbing task which aims to edit the input video through audio, and single image animation using the audio as conditions.
% into two categories based on the type of input visual information:  video-to-video methods and one-shot methods. The visual input of the former methods is a video, and the visual input of the latter methods is a single reference image. 

\subsection{Audio-based Dubbing in Video Editing}

\subsubsection{Arbitrary-subject methods}
% 1. Arbitrary
% Speech2Vid; LipGAN,Wav2Lip; Everybody’s Talkin’; MM22; SyncFaceTalk; 
Arbitrary-subject methods aim at building a general model that does not need to be retrained for different identities.
Speech2Vid~\cite{Chung17b} can re-dub a source video with a different segment of audio thanks to the context encoder. %which help reproduce the background. 
%The generated face can be pasted to original video. 
Reconstructing the lower-half face by inpainting is popular recently~\cite{lipgan, wav2lip, park2022synctalkface}. For example, LipGAN~\cite{lipgan} design a neural network to fill the lower-half face as a pose prior. Wav2Lip~\cite{wav2lip} extends LipGAN using a pre-trained SyncNet as the lip-sync discriminator~\cite{syncnet} to generate accurate lip synchronization. Based on Wav2Lip, SyncTalkFace~\cite{park2022synctalkface} involve the audio-lip memory to store the lip motion features implicitly and retrieve them at inference time. Another category of the methods predicts the intermediate representation first, and then, synthesizes the photo-realistic results by image-to-image translation networks, for example, the facial landmarks~\cite{xie2021towards} and the facial landmarks based on 3D face reconstruction~\cite{song2022everybody}. However, all these methods are struggling to synthesize the high-quality results with editable emotion. 

%Some methods do not reconstruct lip region directly but generate latent representations firstly from audio and then render them to video. 
% propose a two-stage paradigm employs facial landmarks as latent representations to predict lip movements. \cite{song2022everybody} generates 3DMM expression coefficients and use a arbitrary-subject video rendering network to construct photo-realistic video.
%But all the above methods can only handle low-resolution videos. It is difficult to collect a large-scale high-resolution talking face dataset. And it is also really challenging to learn robust audio-visual correlation on small-scale HD dataset. Our method extends the audio-visual correlation learned on low-resolution dataset to high-resolution dataset rather than directly learning on it which cleverly avoids this problem. To the best of our knowledge, our proposed system is the first arbitrary-subject audio-driven visual dubbing method that can handle high-resolution videos.

 
\subsubsection{Personalized methods}
% 2. One person one model
% Audio-driven: SyncObama; Neural Voice Puppetry; Live Speech Portraits;  LipSync3d; AD-Nerf; FACIAL; EVP
% Video-driven: DVP; Neural-style preserved;
% text-driven: write-a-speaker;

% landmark-based: EVP; Live Speech Portraits
% 3d mesh: LipSync3D
% 3dmm: NVP; FACIAL
% nerf: AD-Nerf
%Personalized Visual Dubbing is concerned with generating synchronizing lip motions with arbitrary speech input for one specific target. 
Personalized visual dubbing is easier than the generic one, since these methods are limited to the certain person in the known environment.
For example, SynthesizeObama~\cite{obama} can synthesize the mouth region of Obama by the audio-to-landmark network. Inspired by the face reenactment methods~\cite{dvp, thies2019deferred}, recent visual dubbing methods focus on generating the intermediate representation from audio, and then, rendering the photo-realistic results by the image-to-image translation networks. For example, several works~{\cite{NVP, AudioDVP,zhang2021facial}} focus on the expression coefficient from the audio features and render the photo-realistic results by the image generation networks~\cite{thies2019deferred, dvp, vid2vid}. Facial landmarks~\cite{lu2021live} and edges~\cite{EVP} are also popular choices by projecting the 3D rendered faces since it contains sparser information.
Furthermore, 3D mesh-based~\cite{lipsync3d} and NeRF~\cite{nerf}-based methods~\cite{adnerf} are also powerful.
% \textcolor{red}{LipSync3D~\cite{lipsync3d} normalizes the input face on pose and lighting, and then utilizes 3D mesh as the intermediate representation to generate video.} 
% \textcolor{red}{NeRF~\cite{nerf}-based method~\cite{adnerf} feeds audio into a conditional implicit function to generate a dynamic neural radiance field.}
% Inspired by recently development of Neural Radiance Fields~(NeRF), audio can also be represented as the scene fields. 
Although these methods can synthesize the photo-realistic results, they have relatively limited applications because they need to retrain the model on the specific person and environment.

\nocite{wang2021hfgi}
%the 3D face based methods, AudioDVP~\cite{AudioDVP} distangle audio to emotion space and content space, and then predicts landmark motions. The predicted motion is transferred to the edge map and rendered to video. Live Speech Portraits~\cite{lu2021live} learns facial dynamics and motions include head poses and upper body motions from the projected audio features and use an image-to-image translation network to synthesize photo-realistic renderings in real-time. These two methods both use landmarks as latent representations. 3DMM~\cite{3dmm}-based methods~\cite{NVP, zhang2021facial} firstly generate expression or pose parameters and then render them to photo-realistic videos. 
%NeRF~\cite{nerf}-based method~\cite{adnerf} feeds audio into a conditional implicit function to generate a dynamic neural radiance field and then generate high-fidelity talking-head video using volume rendering. LipSync3D~\cite{lipsync3d} firstly normalize the input face and then utilize 3D mesh as intermediate representation to generate video.

% \cite{obama} uses RNN to predict mouth shapes from audio features. This method requires many hours training data of one person, which is hard to generalize to different targets. NVP~\cite{nvp} generates expression basis from audio and then renders it to video utilizing a person-specific network. \cite{lu2021live} learns facial dynamics and motions include head poses and upper body motions from the projected audio features and use an image-to-image translation network to synthesize photo-realistic renderings in real-time. LipSync3D~\cite{lipsync3d} firstly normalize the input face and then use 3D mesh as intermediate representations to render photo-realistic video. AD-NeRF~\cite{adnerf} feeds audio into a conditional implicit function to generate a dynamic neural radiance field and then generate high-fidelity talking-head video using volume rendering. FACIAL~\cite{zhang2021facial} synthesize the 3D face animation with realistic motions of lips, head poses, and eye blinks with implicit attribute learning. EVP~\cite{AudioDVP} distangles audio to emotion space and content space, and then predicts landmark motions. The predicted motion is transferred to the edge map and rendered to video.

\subsection{Audio-based Single Image Facial Animation}
% one-shot
% X2face; ATVG; DAVS; RNN; Speech-Driven Animation; Flow-guided; PC-AVS; MakeItTalk; Audio2Head,AAAI-22; ijcai22; StyleHEAT
Different from the visual dubbing, single image face animation aims to generate the animation by single driven audio, and it has also been influenced by the video-driven face animation. For example, \cite{song2018talking} generate the motion from audio using the recurrent neural network, \cite{zhou2019talking} disentangle the input to subject-related information and speech-related information by adversarial representation learning. \cite{vougioukas2020realistic, pc-avs} consider the audio as the latent code and drive the face animation by an image generator.
% propose a temporal GAN and three different discriminators to generate lip movements and facial expressions. PC-AVS~\cite{pc-avs} utilize a style-based generator to reconstruct face and the audio servers as the latent code.
The intermediate representation is also a popular choice in this task. ATVG~\cite{chen2019hierarchical} and MakeItTalk~\cite{makeittalk} first generate the facial landmarks from audio, and then, render the video using a landmark-to-video network. Dense flow field is another active research direction~\cite{2203.04036, fomm}. \cite{hdtf} predict the 3DMM coefficients from audio and then transfer these parameters into a flow-based warping network. \cite{wang2021one, wang2021audio2head} borrow the idea from video driven face animation~\cite{fomm}.








